\section{\ENG{SCORM Player \& Backend Runtime Interaction}}
\label{sec:player_backend_interaction}

Η αλληλεπίδραση ανάμεσα στον \ENG{player} και στον \ENG{backend} πυρήνα είναι συχνή, αλλά όχι συμμετρική. Ο \ENG{player} εξαρτάται από το \ENG{engine} για το \ENG{launch context}, για την παραλαβή και καταγραφή των \ENG{commits} και για τον τερματισμό του \ENG{launch}. Το αντίστροφο δεν συμβαίνει: ο \ENG{engine} δεν γνωρίζει τίποτε για την εσωτερική απόδοση του κελύφους \ENG{HTML} ή για τον τοπικό μηχανισμό προσωρινής αποθήκευσης περιεχομένου.

Η βασική διαδρομή προς το \ENG{backend} βρίσκεται στα εσωτερικά \ENG{routes} \texttt{\ENG{/internal/launch/:launchId/commit}} και \texttt{\ENG{/internal/launch/:launchId/terminate}}. Ο \ENG{browser} δεν επικοινωνεί απευθείας με το \ENG{engine}. Καλεί τα εσωτερικά \ENG{endpoints} του \ENG{player}, τα οποία ελέγχουν την ύπαρξη \ENG{token}, επικυρώνουν το \ENG{commit payload} με \texttt{\ENG{zod}} \ENG{schema} και έπειτα προωθούν το αίτημα στο \ENG{engine} μέσω του \texttt{\ENG{RuntimeCommitService}}.

Το μοντέλο αλληλεπίδρασης αποτυπώνεται με σαφήνεια στο \texttt{\ENG{LaunchPageRenderer}}. Η σελίδα που δημιουργείται κρατά \ENG{in-memory state} με \texttt{\ENG{sequence}}, \texttt{\ENG{payload}}, \texttt{\ENG{terminating}} και \texttt{\ENG{terminated}} \ENG{flags}. Κάθε \texttt{\ENG{SetValue()}} ενημερώνει το τοπικό \ENG{payload}. Κάθε \ENG{commit} στέλνει συγκεντρωτικά το τρέχον σύνολο μεταβολών, ενώ υπάρχει και περιοδικός μηχανισμός \ENG{auto-commit} μέσω \texttt{\ENG{setInterval()}}. Κατά το \ENG{beforeunload}, η σελίδα χρησιμοποιεί \texttt{\ENG{navigator.sendBeacon()}} για να αποστείλει ένα τελευταίο \ENG{commit}, γεγονός που δείχνει ότι ο \ENG{player} επιδιώκει ελαχιστοποίηση απώλειας δεδομένων ακόμη και σε μη κανονικούς τερματισμούς του παραθύρου.

Το Σχήμα~\ref{fig:player_backend_interaction} καθιστά σαφή αυτή τη διαδοχή \ENG{browser -> player internal routes -> engine runtime API}. Το σημείο αυτό είναι ουσιώδες για τη συνολική ερμηνεία του συστήματος: ο \ENG{player} λειτουργεί ταυτόχρονα ως κέλυφος παρουσίασης, μηχανισμός διαμεσολάβησης περιεχομένου και διαμεσολαβητής της επικοινωνίας εκτέλεσης, αλλά όχι ως κύριο αποθετήριο της συνεδρίας.

\begin{figure}[H]
  \centering
\begin{chapterfigurebox}
  \begin{tikzpicture}[
      font=\footnotesize,
      node distance=1.8cm and 2.0cm,
      box/.style={draw, rounded corners, fill=gray!5, align=center, text width=3.2cm, minimum height=1.1cm},
      emph/.style={draw, rounded corners, fill=gray!12, align=center, text width=3.5cm, minimum height=1.1cm, font=\footnotesize\bfseries},
      store/.style={draw, cylinder, shape border rotate=90, aspect=0.28, minimum height=1.2cm, minimum width=3.2cm, align=center, fill=gray!5},
      arrow/.style={-Latex, thick}
    ]
    \node[box] (sco) {\ENG{SCO in iframe}};
    \node[box, right=of sco] (host) {\ENG{Player host page}\\\ENG{API bridge}};
    \node[box, right=of host] (internal) {\ENG{Internal player}\\\ENG{routes}};
    \node[emph, below=2.0cm of internal] (engine) {\ENG{Engine runtime API}};
    \node[store, below=1.8cm of engine] (redis) {\ENG{Redis}\\\ENG{active state}};
    \node[store, left=of redis] (postgres) {\ENG{Postgres}\\τελικό αποτέλεσμα};

    \draw[arrow] (sco) -- node[above, align=center, text width=3.1cm]{\ENG{SetValue} / \ENG{Commit} / \ENG{Terminate}} (host);
    \draw[arrow] (host) -- node[above, align=center]{\texttt{\ENG{/internal/}}\\\texttt{\ENG{launch/*}}} (internal);
    \draw[arrow] (internal) -- node[right, align=center]{\ENG{HTTP runtime}\\\ENG{calls}} (engine);
    \draw[arrow] (engine) -- (redis);
    \draw[arrow] (engine) -- node[above, align=center, text width=3.0cm]{συγχρονισμός σε\\\ENG{terminate / timeout}} (postgres);
  \end{tikzpicture}
\end{chapterfigurebox}
  \caption{Αλληλεπίδραση \ENG{player} και \ENG{backend} πυρήνα, από τις κλήσεις \ENG{SCORM API} στην πλευρά του \ENG{browser} έως την ενημέρωση της κατάστασης εκτέλεσης.}
  \label{fig:player_backend_interaction}
\end{figure}

